\chapter{Parkinson Disease}
\index{Parkinson Disease}
\label{sec:Parkinson Disease}

\section{Overview}
Parkinson disease (PD) is a progressive neurodegenerative disorder affecting 1--2\% of adults over 65. Geriatric-specific considerations include Treatment for orthostatic hypotension (common with levodopa and age-related autonomic dysfunction), careful use of dopamine agonists (increased risk of impulse control disorders, hallucinations, and sedation in older adults), and switching to levodopa monotherapy when possible. Deep brain stimulation (DBS) is effective but requires careful patient selection in older adults. Falls are a major concern, and physical therapy (LSVT BIG, balance training) is essential. This neurological disorder affects the central or peripheral nervous system, leading to progressive or episodic symptoms. It significantly impacts quality of life and often requires multidisciplinary care. Neurological disorders affect the central or peripheral nervous system, leading to a wide range of symptoms affecting movement, sensation, cognition, and autonomic function. The nervous system's complexity means that even small disruptions can cause significant functional impairment. Many neurological conditions are chronic and progressive, requiring long-term management and rehabilitation. Advances in neuroimaging and molecular diagnostics have improved understanding and treatment of these conditions. This condition affects many people worldwide and can range from mild to severe. Recognizing the condition early and managing it properly gives you the best chance for a good outcome. Understanding what causes this condition helps your doctor choose the right treatment for you. Learning about your condition and taking an active role in your care are key to managing it long term.

\section{Causes}
Neurological dysfunction can result from structural abnormalities, neurodegenerative processes, demyelination, vascular injury, or neurotransmitter imbalances. Protein aggregation and accumulation is a common mechanism in many neurodegenerative disorders. Excitotoxicity, oxidative stress, and neuroinflammation contribute to neuronal injury. Genetic mutations account for some neurological conditions, while others are sporadic or environmentally triggered. This condition develops from a combination of genetic and environmental factors that disrupt normal body processes. Several factors can work together to trigger or make the condition worse.

\section{Risk Factors}

\begin{itemize}[noitemsep]
  \item Genetic predisposition and family history
  \item Advancing age
  \item Lifestyle factors (diet, exercise, smoking, alcohol)
  \item Environmental exposures
  \item Underlying medical conditions
  \item Medication use
\end{itemize}

\section{Signs and Symptoms}

\begin{itemize}[noitemsep]
  \item Pain or discomfort in the affected area
  \item Difficulty performing daily activities
  \item Changes in how the affected area looks or works
  \item General symptoms may include fatigue, fever, or weight changes
  \item Symptoms may develop slowly or appear suddenly
\end{itemize}

\section{Progression and Stages}
This condition can progress through different stages over time. Early stages often involve mild symptoms that you may not notice right away. As the condition progresses, symptoms become more noticeable and may affect your daily activities. Advanced stages may involve more serious symptoms and complications.

\section{Complications}
In the geriatric population, PD is linked to higher rates of postural instability, falls, dementia, difficulty swallowing, and medication complications (motor fluctuations, dyskinesias, psychosis).

\begin{itemize}[noitemsep]
  \item The condition may worsen over time if not treated
  \item Increased risk of other infections or injuries
  \item Side effects from medications or other treatments
  \item Reduced quality of life
  \item Emotional and psychological effects
\end{itemize}

\section{Diagnosis}

\begin{itemize}[noitemsep]
  \item Your doctor will take a detailed medical history and perform a physical exam
  \item Lab tests to check how your organs are working
  \item Imaging tests (like X-rays or MRI) to look at the affected area
  \item Specialized tests depending on which body system is involved
  \item Referral to a specialist for further evaluation
\end{itemize}

\section{Prevention}

\begin{itemize}[noitemsep]
  \item Eat a balanced diet and exercise regularly
  \item Avoid known triggers and risk factors
  \item Go for regular check-ups so any problems are caught early
  \item Follow your doctor's screening recommendations
\end{itemize}

\section{Treatment Options}

\begin{itemize}[noitemsep]
  \item Medications to control symptoms and treat the underlying cause
  \item Lifestyle changes and self-care
  \item Physical therapy and rehabilitation
  \item Surgery when needed
  \item Regular follow-up care
\end{itemize}

\section{Living with It}

\begin{itemize}[noitemsep]
  \item Follow your treatment plan as prescribed
  \item Keep up with regular appointments with your healthcare provider
  \item Adjust your daily habits as needed
  \item Reach out to family, friends, or support groups for help
  \item Keep learning about your condition and how to manage it
\end{itemize}

\section{Outlook}
Outlook is progressive. Neurological conditions vary widely in their course and outcomes. Some are progressive, while others follow a relapsing-remitting pattern. Early intervention and rehabilitation can improve outcomes. Neurological conditions vary significantly in their course, from acute and self-limited to chronic and progressive. Early diagnosis and intervention can slow progression and improve quality of life. Rehabilitation and supportive therapies play crucial roles in maintaining function. Research continues to develop disease-modifying treatments for previously untreatable conditions. Multidisciplinary care addressing medical, functional, and psychosocial needs optimizes outcomes.
